\section{The Hierarchy of One Size}\label{sec:hierarchy}

Fix $s$.
Two nuclei of size $s$ are disjoint or nested: if $N$ is an $\Nuc{s}{k}$ and $N'$ an $\Nuc{s}{k'}$ with $k\le k'$ and $N\cap N'\ne\emptyset$, then $N'\subseteq N$, because $N\cup N'$ satisfies the level-$k$ condition and $N$ is maximal.
The nuclei of size $s$ therefore form a laminar family, and the same vertex set can be a nucleus for a whole interval of levels.

\begin{definition}[canonical node, own node]\label{def:node}
A \emph{canonical node} of size $s$ is a vertex set that is an $\Nuc{s}{k}$ for at least one $k$; its \emph{top} is the largest such $k$, which equals the smallest core value of its vertices.
The canonical nodes of size $s$ ordered by inclusion form the forest $\Tree{s}$.
The \emph{own node} $\own{s}{v}$ of a vertex with $\core{s}{v}\ge1$ is the $\Nuc{s}{\core{s}{v}}$ containing $v$: the deepest node of $\Tree{s}$ containing $v$, whose top is $\core{s}{v}$.
\end{definition}

For $1\le k\le\core{s}{v}$ the community of $v$ at $(s,k)$ is the unique ancestor-or-self of $\own{s}{v}$ in $\Tree{s}$ whose level interval contains $k$; equivalently, the highest ancestor whose top is at least $k$, since tops decrease strictly along every root path.
In the running example $\Tree{2}$ is a root of top 3 holding all 15 vertices with two children of top 4, $A$ and $W\cup\{x\}$; $\Tree{3}$ has two roots, $N=B\cup W\cup\{x\}$ of top 3 (the edge $a_5b_3$ lies in no triangle, so $A$ is not joined to it) and $A$ of top 6; $\Tree{4}$ has the roots $A$ (top 4) and $B$ (top 1); $\Tree{5}$ is $A$ alone.
The own node of $b_1$ at size 2 is the root, at size 3 it is $N$, at size 4 it is $B$.

\begin{fact}[support nesting]\label{fact:nesting}
$\core{s}{v}\ge1$ if and only if $2\le s\le\omg{v}$.
\end{fact}

\begin{theorem}[floor and certified tail]\label{thm:tail}
Let $2\le s\le\omg{v}$ and $f_s(v)=\binomc{\omg{v}-1}{s-1}$.
\begin{enumerate}[label=(\alph*),leftmargin=*,nosep]
\item $\core{s}{v}\ge f_s(v)$.
\item If $\core{s}{v}=f_s(v)$ and $s<\omg{v}$, then $\core{s+1}{v}=f_{s+1}(v)$.
\end{enumerate}
Hence once the value of $v$ meets its floor it stays on the floor at every larger size.
\end{theorem}
\begin{proof}
(a) A largest clique $M$ containing $v$ has $\omg{v}$ vertices; every vertex of $M$ lies in $\binomc{\omg{v}-1}{s-1}$ $s$-cliques of $M$, and $M$ is connected through them, so $M$ lies in an $\Nuc{s}{f_s(v)}$.

(b) We use the Kruskal--Katona theorem in Lov\'asz's form \cite{kruskal1963,katona1968,lovasz1979}: $k\ge1$ distinct $s$-subsets of a ground set have at least $g_s(k)=\binomc{x}{s-1}$ distinct $(s-1)$-subsets, where $x\ge s$ is the real number with $\binomc{x}{s}=k$; both $x\mapsto\binomc{x}{s}$ and $x\mapsto\binomc{x}{s-1}$ increase strictly for $x\ge s-1$.
Let $N$ be the $\Nuc{s+1}{\core{s+1}{v}}$ containing $v$ and take any $u\in N$.
Removing $u$ from the at least $\core{s+1}{v}$ $(s+1)$-cliques of $N$ through $u$ leaves as many distinct $s$-subsets of the neighbours of $u$ in $N$; their $(s-1)$-shadow has at least $g_s(\core{s+1}{v})$ members, and each member together with $u$ is an $s$-clique of $N$ through $u$.
So every vertex of $N$ lies in at least $g_s(\core{s+1}{v})$ $s$-cliques of $N$, and $N$ is connected through them (two $(s+1)$-cliques that share a vertex contain $s$-cliques that share it); therefore $\core{s}{v}\ge g_s(\core{s+1}{v})$.
Now suppose $\core{s+1}{v}>f_{s+1}(v)=\binomc{\omg{v}-1}{s}$.
Then the $x$ with $\binomc{x}{s}=\core{s+1}{v}$ exceeds $\omg{v}-1\ge s$, so $\core{s}{v}\ge g_s(\core{s+1}{v})=\binomc{x}{s-1}>\binomc{\omg{v}-1}{s-1}=f_s(v)$, contradicting $\core{s}{v}=f_s(v)$.
Since $\core{s+1}{v}\ge f_{s+1}(v)$ by (a), equality holds.
\end{proof}

We write $\sig{v}$ for the smallest $s$ at which equality holds ($\sig{v}\le\omg{v}+1$, the latter meaning that it never holds).
In the example the vertices of $A$ and of $B$ are certified from size 2 ($\kappa_2(a_i)=4=\binomc{4}{1}$, $\kappa_2(b_i)=3=\binomc{3}{1}$), $x$ from size 3 ($\kappa_3(x)=3=\binomc{3}{2}$ while $\kappa_2(x)=4>3$), and the vertices of $W$ never ($\sig{u_i}=4=\omg{u_i}+1$).
For $s\ge\sig{v}$ the value $\core{s}{v}$ is $\binomc{\omg{v}-1}{s-1}$ and need not be stored; the values for $2\le s<\sig{v}$ are the \emph{residue} of $v$.
On real graphs most values are in the tail: \dblp has 1,246,933 pairs $(v,s)$ with $\core{s}{v}\ge1$ and 84,010 residue values.
